\chapter{Esercizi: Ereditarietà e polimorfismo in Java}
\label{chap:esercizi-ereditarieta-polimorfismo-java}

\begin{exercise}
\label{ex:fi-esercizi-java-ereditarieta-ordinamento-confronto}
Perché il confronto tra oggetti richiede un criterio definito dal loro tipo?
\end{exercise}

A differenza dell’ordine numerico, l’ordine tra oggetti non è universale:
dipende dalle proprietà scelte dalla classe. Solo oggetti confrontabili
secondo un criterio possono essere ordinati
(\cref{sec:fi-java-ereditarieta-ordinamento-oggetti}).

\begin{exercise}
\label{ex:fi-esercizi-java-ereditarieta-compareto}
Quale significato ha il risultato di \texttt{String.compareTo()}?
\end{exercise}

Un valore negativo indica che la prima stringa precede la seconda, zero che
risultano equivalenti nel confronto e un valore positivo che la segue
(\cref{sec:fi-java-ereditarieta-ordinamento-oggetti}).

\begin{exercise}
\label{ex:fi-esercizi-java-ereditarieta-selection-stringhe}
Come cambia \texttt{selectionSort()} passando da un array di interi a uno di
stringhe?
\end{exercise}

Si usa un array di \texttt{String} e si sostituisce il confronto numerico
con \texttt{compareTo()} a risultato negativo. La logica di ricerca del
minimo non cambia (\cref{sec:fi-java-ereditarieta-ordinamento-oggetti}).

\begin{exercise}
\label{ex:fi-esercizi-java-ereditarieta-subclass}
Che cosa significa che \texttt{SavingsAccount} è sottoclasse di
\texttt{BankAccount}?
\end{exercise}

Ogni istanza di \texttt{SavingsAccount} è anche un \texttt{BankAccount}; ne
riutilizza i comportamenti accessibili e può introdurre membri specifici
(\cref{def:fi-java-ereditarieta}).

\begin{exercise}
\label{ex:fi-esercizi-java-ereditarieta-eredita}
Che cosa può ereditare una sottoclasse?
\end{exercise}

Può utilizzare i membri accessibili della superclasse e sovrascrivere i
metodi di istanza che lo consentono. I campi \texttt{private} non sono
accessibili direttamente dalla sottoclasse
(\cref{sec:fi-java-ereditarieta-conti}).

\begin{exercise}
\label{ex:fi-esercizi-java-ereditarieta-object}
Perché \texttt{Object} è detta superclasse universale?
\end{exercise}

Perché è la radice della gerarchia delle classi Java: una classe che non
dichiara una superclasse deriva implicitamente da \texttt{Object}
(\cref{def:fi-java-ereditarieta-object}).

\begin{exercise}
\label{ex:fi-esercizi-java-ereditarieta-duplicazione}
Perché non conviene copiare il codice di \texttt{BankAccount} in
\texttt{SavingsAccount}?
\end{exercise}

Si duplicano comportamento e responsabilità di manutenzione. L’ereditarietà
permette alla sottoclasse di utilizzare la stessa implementazione, senza
replicarla (\cref{sec:fi-java-ereditarieta-conti}).

\begin{exercise}
\label{ex:fi-esercizi-java-ereditarieta-addinterest}
Perché \texttt{addInterest()} invoca \texttt{getBalance()} e
\texttt{deposit()} invece di modificare \texttt{balance}?
\end{exercise}

Il campo \texttt{balance} è \texttt{private} in \texttt{BankAccount}. I
metodi pubblici permettono alla sottoclasse di operare sul saldo
rispettandone l’incapsulamento (\cref{sec:fi-java-ereditarieta-conti}).

\begin{exercise}
\label{ex:fi-esercizi-java-ereditarieta-metodi-tipi}
Che differenza c’è tra un metodo nuovo, ereditato e sovrascritto?
\end{exercise}

Un metodo nuovo compare soltanto nella sottoclasse; uno ereditato conserva
l’implementazione della superclasse; uno sovrascritto ne fornisce una nuova
implementazione (\cref{sec:fi-java-ereditarieta-metodi}).

\begin{exercise}
\label{ex:fi-esercizi-java-ereditarieta-override-firma}
Quali condizioni deve rispettare un metodo per sovrascriverne uno della
superclasse?
\end{exercise}

Deve essere un metodo di istanza sovrascrivibile, con firma corrispondente,
tipo di ritorno compatibile e accessibilità non più restrittiva.
\texttt{@Override} consente al compilatore di verificare la sovrascrittura
(\cref{def:fi-java-ereditarieta-override}).

\begin{exercise}
\label{ex:fi-esercizi-java-ereditarieta-this-super}
Qual è la differenza tra \texttt{this} e \texttt{super}?
\end{exercise}

\texttt{this} denota l'istanza corrente; \texttt{super} permette di
selezionare un membro della superclasse accessibile alla sottoclasse,
per esempio l'implementazione di un metodo sovrascritto. Entrambi
operano sulla stessa istanza (\cref{sec:fi-java-ereditarieta-metodi}).

\begin{exercise}
\label{ex:fi-esercizi-java-ereditarieta-upcast}
Perché una variabile di tipo \texttt{BankAccount} può indicare un oggetto
\texttt{SavingsAccount}?
\end{exercise}

Perché il tipo della sottoclasse è compatibile con quello della superclasse.
L’upcasting copia soltanto il riferimento; l’oggetto non viene convertito né
duplicato (\cref{sec:fi-java-ereditarieta-conversioni}).

\begin{exercise}
\label{ex:fi-esercizi-java-ereditarieta-instanceof}
A che cosa serve \texttt{instanceof}?
\end{exercise}

Verifica durante l’esecuzione se l’oggetto indicato è compatibile con il
tipo specificato, permettendo fra l’altro di proteggere un cast verso una
sottoclasse (\cref{sec:fi-java-ereditarieta-instanceof}).

\begin{exercise}
\label{ex:fi-esercizi-java-ereditarieta-instanceof-tipo}
Il risultato di \texttt{instanceof} dipende dal tipo dichiarato o dal tipo
effettivo dell’oggetto?
\end{exercise}

Dipende dal tipo effettivo dell’oggetto; applicato a \texttt{null}
restituisce \texttt{false} (\cref{sec:fi-java-ereditarieta-instanceof}).

\begin{exercise}
\label{ex:fi-esercizi-java-ereditarieta-polimorfismo}
Che cosa significa polimorfismo per sottotipi in Java?
\end{exercise}

Un riferimento al tipo della superclasse può indicare istanze di sottoclassi
diverse; l’esecuzione di un metodo di istanza sovrascritto seleziona
l’implementazione in base al tipo effettivo dell’oggetto
(\cref{def:fi-java-ereditarieta-polimorfismo}).

\begin{exercise}
\label{ex:fi-esercizi-java-ereditarieta-binding}
Come si distinguono selezione anticipata e selezione dinamica dei metodi?
\end{exercise}

La risoluzione dei sovraccarichi usa i tipi statici in compilazione
(\emph{early binding}); per i metodi di istanza sovrascritti,
l'implementazione viene scelta in base al tipo effettivo dell'oggetto
durante l'esecuzione (\emph{late binding})
(\cref{sec:fi-java-ereditarieta-polimorfismo}).

\begin{exercise}
\label{ex:fi-esercizi-java-ereditarieta-cast-sicuro}
Scrivi un esempio in cui \texttt{instanceof} protegge il cast verso
\texttt{SavingsAccount}.
\end{exercise}

\begin{unipdcode}[language=Java]
public static void accreditaInteressi(BankAccount conto) {
    if (conto instanceof SavingsAccount) {
        SavingsAccount s = (SavingsAccount) conto;
        s.addInterest();
    }
}
\end{unipdcode}
Il metodo, da dichiarare in una classe, accetta anche un conto che non
sia di risparmio: il cast viene eseguito soltanto quando l'oggetto è
compatibile (\cref{sec:fi-java-ereditarieta-instanceof}).
